% Options for packages loaded elsewhere
\PassOptionsToPackage{unicode}{hyperref}
\PassOptionsToPackage{hyphens}{url}
\documentclass[
  ignorenonframetext,
]{beamer}
\newif\ifbibliography
\usepackage{pgfpages}
\setbeamertemplate{caption}[numbered]
\setbeamertemplate{caption label separator}{: }
\setbeamercolor{caption name}{fg=normal text.fg}
\beamertemplatenavigationsymbolsempty
% remove section numbering
\setbeamertemplate{part page}{
  \centering
  \begin{beamercolorbox}[sep=16pt,center]{part title}
    \usebeamerfont{part title}\insertpart\par
  \end{beamercolorbox}
}
\setbeamertemplate{section page}{
  \centering
  \begin{beamercolorbox}[sep=12pt,center]{section title}
    \usebeamerfont{section title}\insertsection\par
  \end{beamercolorbox}
}
\setbeamertemplate{subsection page}{
  \centering
  \begin{beamercolorbox}[sep=8pt,center]{subsection title}
    \usebeamerfont{subsection title}\insertsubsection\par
  \end{beamercolorbox}
}
% Prevent slide breaks in the middle of a paragraph
\widowpenalties 1 10000
\raggedbottom
\AtBeginPart{
  \frame{\partpage}
}
\AtBeginSection{
  \ifbibliography
  \else
    \frame{\sectionpage}
  \fi
}
\AtBeginSubsection{
  \frame{\subsectionpage}
}
\usepackage{iftex}
\ifPDFTeX
  \usepackage[T1]{fontenc}
  \usepackage[utf8]{inputenc}
  \usepackage{textcomp} % provide euro and other symbols
\else % if luatex or xetex
  \usepackage{unicode-math} % this also loads fontspec
  \defaultfontfeatures{Scale=MatchLowercase}
  \defaultfontfeatures[\rmfamily]{Ligatures=TeX,Scale=1}
\fi
\usepackage{lmodern}
\ifPDFTeX\else
  % xetex/luatex font selection
\fi
% Use upquote if available, for straight quotes in verbatim environments
\IfFileExists{upquote.sty}{\usepackage{upquote}}{}
\IfFileExists{microtype.sty}{% use microtype if available
  \usepackage[]{microtype}
  \UseMicrotypeSet[protrusion]{basicmath} % disable protrusion for tt fonts
}{}
\makeatletter
\@ifundefined{KOMAClassName}{% if non-KOMA class
  \IfFileExists{parskip.sty}{%
    \usepackage{parskip}
  }{% else
    \setlength{\parindent}{0pt}
    \setlength{\parskip}{6pt plus 2pt minus 1pt}}
}{% if KOMA class
  \KOMAoptions{parskip=half}}
\makeatother
\usepackage{longtable,booktabs,array}
\newcounter{none} % for unnumbered tables
\usepackage{calc} % for calculating minipage widths
\usepackage{caption}
% Make caption package work with longtable
\makeatletter
\def\fnum@table{\tablename~\thetable}
\makeatother
\setlength{\emergencystretch}{3em} % prevent overfull lines
\providecommand{\tightlist}{%
  \setlength{\itemsep}{0pt}\setlength{\parskip}{0pt}}
\usepackage{bookmark}
\IfFileExists{xurl.sty}{\usepackage{xurl}}{} % add URL line breaks if available
\urlstyle{same}
\hypersetup{
  hidelinks,
  pdfcreator={LaTeX via pandoc}}

\author{\texorpdfstring{}{}}
\date{}

\begin{document}

\begin{frame}[fragile]{The Evidence: What We Proved in Part 2}
\protect\phantomsection\label{sec:paper2-review}
Part 1 was the theory. Part 2 was the proof. We ran \textbf{1,557 tests}
against \textbf{950 attack samples} across \textbf{6 production
architectures}.

Here is what the data says.

\begin{block}{The Experimental Setup}
\protect\phantomsection\label{the-experimental-setup}
We tested the six most common multi-agent architectures found in
production:

\begin{itemize}
\tightlist
\item
  \textbf{Claude Code} (Hierarchical)
\item
  \textbf{AutoGPT} (Autonomous Loop)
\item
  \textbf{CrewAI} (Role-Based Team)
\item
  \textbf{LangGraph} (State Machine)
\item
  \textbf{MetaGPT} (SOP-Driven)
\item
  \textbf{Camel} (Peer-to-Peer)
\end{itemize}

The test corpus included direct prompt injection, poisoned RAG contexts,
deep trust exploitation, and multi-turn social engineering.
\end{block}

\begin{block}{Finding 1: Defense Layering vs.~Individual Efficacy}
\protect\phantomsection\label{finding-1-defense-layering-vs.-individual-efficacy}
\textbf{The Data}: Individual defenses (like just a firewall) stopped
\textasciitilde60-70\% of attacks. The full CIF stack stopped
\textbf{97\%}. \textbf{The Implication}: The defenses demonstrated
orthogonal coverage. The firewall blocked inputs that the sandbox would
have missed, and the sandbox identified anomalies that the trust
calculus would have permitted. The data suggests that removing any
single layer creates a statistically significant vulnerability gap.
\end{block}

\begin{block}{Finding 2: State Machine Determinism}
\protect\phantomsection\label{finding-2-state-machine-determinism}
\textbf{The Data}: \textbf{LangGraph} architectures achieved the highest
detection rates (98\%). \textbf{The Mechanism}: In a state machine,
valid transitions are explicitly defined. Attempts to move from an
\texttt{Analyze} state to an \texttt{Execute} state without passing
through \texttt{Verify} were caught immediately. \textbf{The
Implication}: Explicit state definitions provide a ``security for free''
effect, where the architecture itself enforces behavioral invariants
that are difficult to bypass.
\end{block}

\begin{block}{Finding 3: Hierarchical Vulnerability}
\protect\phantomsection\label{finding-3-hierarchical-vulnerability}
\textbf{The Data}: \textbf{Claude Code} and \textbf{AutoGPT} styles were
strong against external attacks but vulnerable to \textbf{Systemic
(\(\Omega_5\))} compromise. \textbf{The Mechanism}: When the ``Boss''
agent was tricked, it ordered ``Worker'' agents to execute malicious
actions, and they complied. \textbf{The Implication}: Hierarchical
systems exhibit a Single Point of Failure at the root. Security in these
topologies requires worker-level tripwires that can reject orders even
from authenticated superiors.
\end{block}

\begin{block}{Finding 4: Roles as Security Boundaries}
\protect\phantomsection\label{finding-4-roles-as-security-boundaries}
\textbf{The Data}: \textbf{CrewAI} performed exceptionally well against
privilege escalation. \textbf{The Mechanism}: ``Roles'' acted as
effective containers for capabilities. A ``Researcher'' agent lacked the
tools to write code, rendering code-execution attacks against it inert.
\textbf{The Implication}: Strict role definitions function as effective
security boundaries, limiting the blast radius of a compromised agent.
\end{block}

\begin{block}{Finding 5: The Stealth-Impact Tradeoff}
\protect\phantomsection\label{finding-5-the-stealth-impact-tradeoff}
\textbf{The Data}: High-impact attacks were consistently easier to
detect than low-impact nudges. \textbf{The Implication}:
``Catastrophic'' takeover attempts generate significant noise in the
system state. The primary challenge for defenders is not the sudden
takeover, but the slow, subtle drift of agent beliefs over long time
horizons.
\end{block}
\end{frame}

\begin{frame}
\begin{block}{A Note on the Numbers}
\protect\phantomsection\label{a-note-on-the-numbers}
The detection rates in Part 2 are derived from a calibrated parametric
simulation, modeled on the architecture's topology. They represent the
\emph{structural} security of the design.

\begin{itemize}
\tightlist
\item
  \textbf{97\% Detection} means: ``In a fully implemented system of this
  topology, 97\% of these attack vectors violate a defined constraint.''
\item
  It does \textbf{not} mean: ``We have a magic Python script that
  catches 97\% of all evil AI thoughts.''
\end{itemize}

We proved the \emph{architecture} works. The implementation fidelity is
the variable for the builder.

\begin{block}{Tripwire Configuration Data}
\protect\phantomsection\label{tripwire-configuration-data}
The simulations utilized the following tripwire densities to achieve the
reported results:

{\def\LTcaptype{none} % do not increment counter
\begin{longtable}[]{@{}lll@{}}
\toprule\noalign{}
Category & Count Used in Sim & Placement Strategy \\
\midrule\noalign{}
\endhead
Identity canaries & 3+ per agent & Core identity beliefs \\
Boundary canaries & 5+ per agent & Permission boundaries \\
Principal canaries & 2+ per agent & Trust relationships \\
Temporal canaries & 1 per agent & Session continuity \\
\bottomrule\noalign{}
\end{longtable}
}
\end{block}
\end{block}
\end{frame}

\end{document}
